This metronome plugin generates a rhythmic click track with
customizable timing, beat signatures, and filtered sound
characteristics. It supports synchronization to the transport timeline
or free-running operation, allowing real-time control of parameters
like BPM and volume via OSC. Additionally, it features a dispatch
mechanism that allows users to schedule and automatically send custom
OSC messages on specific downbeats, enabling it to trigger events in
other parts of the session.

\input{tabmetronome.tex}

%OSC messages can be dispatched on beat one using the ``/dispatchin'' OSC variables.


\definecolor{shadecolor}{RGB}{236,236,255}\begin{snugshade}
{\footnotesize
\label{osctab:tascarapmetronome}
OSC variables:
\nopagebreak

\begin{tabularx}{\textwidth}{llllX}
\hline
path & fmt. & range & readable & description\\
\hline
\attr{/.../a1} & f &  & yes & Amplitude of the first beat (the "downbeat").\\
\attr{/.../ao} & f &  & yes & Amplitude of the other beats.\\
\attr{/.../bpb} & i &  & no & Beats per bar. This is a vector allowing for varying bar lengths (e.g., "4 3" for a 4/4 bar followed by a 3/4 bar).\\
\attr{/.../bpm} & f &  & yes & Beats per minute. Defaults to 120.\\
\attr{/.../bypass} & i & bool & yes & If true, the metronome sound is muted.\\
\attr{/.../changeonone} & i & bool & yes & If true, parameter changes (like BPM or volume updates) are applied only on the next "one" (the first beat of a bar). If false, changes are applied immediately.\\
\attr{/.../dispatchin} & i &  & yes & A counter used to trigger an OSC message dispatch on a specific beat. (See detailed explanation below).\\
\attr{/.../dispatchmsg} & (any) &  & no & An OSC method that accepts a generic message. The data sent here is stored and used later by the dispatch mechanism.\\
\attr{/.../dispatchpath} & s & string & yes & The OSC path (address) that the stored message will be sent to when the dispatch is triggered.\\
\attr{/.../filter/f1} & f &  & yes & Resonance frequency of the first beat.\\
\attr{/.../filter/fo} & f &  & yes & Resonance frequency of other beats.\\
\attr{/.../filter/q1} & f &  & yes & Filter resonance (Q factor) of the first beat.\\
\attr{/.../filter/qo} & f &  & yes & Filter resonance (Q factor) of other beats.\\
\attr{/.../sync} & i & bool & yes & Enables synchronization to the transport time line. If false (default), the metronome runs freely based on the audio sample rate.\\
\hline
\end{tabularx}
}
\end{snugshade}
\definecolor{shadecolor}{RGB}{255,230,204}


The \verb!dispatchin! feature allows the metronome to act as a
sequencer or trigger for other OSC events. It enables you to schedule
an OSC message to be sent out automatically when the metronome reaches
a specific downbeat. How it works:

\begin{enumerate} 
\item {\bf Setup (The Payload):} You first define the message you want to send. You send an OSC message to the \verb!/dispatchmsg! method of this plugin containing the arguments (e.g., floats, strings, integers) you want to dispatch. Simultaneously, you set the \verb!/dispatchpath! variable to the target OSC address (e.g., \verb!/scene/start!) where this message should eventually be sent.
\item {\bf Scheduling (The Counter):} You set the variable \verb!/dispatchin! to an integer value. This value represents the number of bars to wait before the message is sent.
\begin{itemize}
\item If you set \verb!/dispatchin! to 1, the message is dispatched immediately on the very next downbeat.
\item If you set \verb!/dispatchin! to 4, the plugin will count down 3, 2, 1 during the subsequent downbeats, and dispatch the message on the 4th downbeat.
\end{itemize}
\item {\bf Execution:} Inside the audio processing loop (\verb!ap_process!), the plugin checks every time it hits a "one" (the first beat of a bar, where \verb!beat! == 0):
\begin{itemize}
\item It checks if dispatchin is greater than 0.
\item If so, it decrements the counter (\verb!--dispatchin!).
\item If the counter reaches exactly 0 during this beat, the plugin calls \verb!srv_->dispatch_data_message!. This takes the stored message payload and sends it out to the stored \verb!/dispatchpath!.
\end{itemize}
\end{enumerate}


Each sub-message can be defined using a \elem{msg} element.

\input{tabmsg.tex}
